% CS 675 Machine Learning, Fall 2026, section 1B1. Three short slides for a
% live session: a question about hyperparameters (blank below, for the
% discussion); one-vs-rest units side by side against a softmax layer, with
% the link to the Newark seasons demo where the two models disagree; and
% nearest-neighbour regression written as softmax attention, with a worked
% picture from eight Newark days. The numbers of the third slide come from
% attention-data.tex, written by prep.py; the figure of the second is the
% CS 301 figure seasons-two-models.pdf (fig-networks.py in that repo).
% Build: pdflatex three-pictures.tex  (run twice)
\DocumentMetadata{lang=en-US,pdfversion=1.7,tagging=on}
\documentclass[aspect-ratio=16:9,11pt,frame-title-arg]{ltx-talk}
\usepackage{cs675-talk}
% written by attention_tex.py from attention.json; do not edit
\def\AttQuery{71}
\def\AttTau{8}
\def\AttYhat{49.5}
\def\AttKnn{53.6}
\def\AttLine{53.4}
\def\AttDays{32.0/25.2/0.000/2017-01-31, 41.0/27.1/0.000/2013-11-29, 50.0/30.9/0.014/2013-03-05, 57.0/43.0/0.098/2016-10-31, 66.0/46.9/0.373/2013-04-16, 75.0/50.0/0.400/2015-10-06, 84.9/64.0/0.100/2015-09-18, 91.9/66.9/0.015/2008-07-03}


% ---------- the two multiclass pictures ----------
\tikzset{
  inp/.style={circle,fill=ink,inner sep=0pt,minimum size=1.8mm},
  unit/.style={circle,draw=ink,fill=paper,line width=0.6pt,minimum size=5.2mm,inner sep=0pt,font=\scriptsize},
  act/.style={circle,draw=njred,fill=njred!12,line width=0.6pt,minimum size=4.6mm,inner sep=0pt,font=\scriptsize},
  lossbox/.style={draw=muted,rounded corners=1pt,fill=white,font=\tiny,inner sep=1.6pt,minimum width=9mm,minimum height=3.6mm},
  lbl/.style={font=\scriptsize,color=muted},
  wire/.style={line width=0.35pt,color=faint},
  arr/.style={-{Latex[length=1.4mm]},line width=0.45pt,color=muted}}

\newcommand{\ovrpic}{%
\begin{tikzpicture}[x=1mm,y=1mm]
  \node[inp] (x1) at (0,19) {}; \node[inp] (x2) at (0,9) {};
  \node[lbl,anchor=east] at (-1.5,19) {$x_1$}; \node[lbl,anchor=east] at (-1.5,9) {$x_2$};
  \foreach \k/\yy in {1/26,2/18,3/10,4/2} {
    \node[unit] (u\k) at (20,\yy) {$s_\k$};
    \draw[wire] (x1) -- (u\k); \draw[wire] (x2) -- (u\k);
    \node[act] (p\k) at (35,\yy) {$\sigma$};
    \draw[arr] (u\k) -- (p\k);
    \node[lossbox] (l\k) at (52,\yy) {loss$_\k$};
    \draw[arr] (p\k) -- (l\k);
  }
  \node[lbl,anchor=south] at (20,30.5) {score};
  \node[lbl,anchor=south] at (35,30.5) {$\hat y_k=\sigma(s_k)$};
  \node[lbl,anchor=south] at (52,30.5) {four losses};
  \node[lbl,anchor=north,align=center] at (35,-2.5) {trained \alert{one at a time}, each on ``is it class $k$?''\\
    outputs need not add to 1; predict by argmax};
\end{tikzpicture}}

\newcommand{\softmaxpic}{%
\begin{tikzpicture}[x=1mm,y=1mm]
  \node[inp] (x1) at (0,19) {}; \node[inp] (x2) at (0,9) {};
  \node[lbl,anchor=east] at (-1.5,19) {$x_1$}; \node[lbl,anchor=east] at (-1.5,9) {$x_2$};
  \foreach \k/\yy in {1/26,2/18,3/10,4/2} {
    \node[unit] (u\k) at (20,\yy) {$s_\k$};
    \draw[wire] (x1) -- (u\k); \draw[wire] (x2) -- (u\k);
  }
  \node[draw=njred,fill=njred!8,rounded corners=1.2pt,line width=0.6pt,minimum width=6mm,minimum height=32mm,inner sep=0pt] (sm) at (34,14) {};
  \node[font=\scriptsize,color=njred,rotate=90] at (34,14) {softmax};
  \foreach \k/\yy in {1/26,2/18,3/10,4/2} {
    \draw[arr] (u\k) -- (31,\yy);
    \node[act] (p\k) at (46,\yy) {$p_\k$};
    \draw[arr] (37,\yy) -- (p\k);
  }
  \node[lossbox,minimum height=7mm] (l) at (60,14) {one loss};
  \foreach \k in {1,2,3,4} { \draw[arr] (p\k) -- (l); }
  \node[lbl,anchor=south] at (20,30.5) {score};
  \node[lbl,anchor=south] at (46,30.5) {$p_k=\dfrac{e^{s_k}}{\sum_j e^{s_j}}$};
  \node[lbl,anchor=north,align=center] at (35,-2.5) {trained \alert{together} on the cross-entropy $-\ln p_{\text{true}}$\\
    outputs add to 1; predict by argmax};
\end{tikzpicture}}

\begin{document}
\deckcard{Three pictures}{hyperparameters \,--\, one-vs-rest against softmax \,--\, nearest neighbours as attention}

% ---------------------------------------------------------------- 1
\questionslide{Hyperparameters}{\bigq{What is a hyperparameter?}}

% ---------------------------------------------------------------- 2
\begin{frame}[t]{Four units side by side, or one softmax layer}
  \vspace{1mm}
  \begin{columns}[c]
    \begin{column}{0.49\textwidth}
      \centering
      {\small\alert{One-vs-rest}: $k$ independent units}\par\vspace{0.5mm}
      \figalt{Two inputs feed four separate units. Each unit's score goes through its own sigmoid to its own loss box, so there are four losses. Caption: trained one at a time, each on the question is it class k; outputs need not add to 1; predict by argmax.}
      \scalebox{0.8}{\ovrpic}
    \end{column}
    \begin{column}{0.49\textwidth}
      \centering
      {\small\alert{Softmax}: $k$ units linked by one layer}\par\vspace{0.5mm}
      \figalt{The same two inputs feed four units, whose four scores enter one softmax box together. The four outputs p one to p four all feed a single loss box. Caption: trained together on the cross-entropy; outputs add to 1; predict by argmax.}
      \scalebox{0.8}{\softmaxpic}
    \end{column}
  \end{columns}
  \vspace{2.5mm}
  \begin{columns}[c]
    \begin{column}{0.46\textwidth}
      \ttab{\scriptsize\setlength{\tabcolsep}{4pt}\renewcommand{\arraystretch}{1.1}
      \begin{tabular}{l|ll}
        & separate units & softmax layer \\ \hline
        output & $\sigma(s_k)$ & $p_k$ \\
        computed from & its own score & all $k$ scores \\
        loss & one per unit & one for the model \\
        trained & separately & jointly \\
      \end{tabular}}
    \end{column}
    \begin{column}{0.52\textwidth}
      {\small Same inputs, same units, \alert{different weights}: on the Newark seasons the two models name a different season on 160 of 1{,}457 validation days.}
      \par\vspace{1mm}
      {\scriptsize\muteds{Notes: 5.5 (one-vs-rest), 16.2 and 16.3 (softmax, cross-entropy).}}
    \end{column}
  \end{columns}
  \vspace{1.5mm}
  {\footnotesize\muteds{A day the two disagree on:}\ \scriptsize\url{https://ikoutis.github.io/courses/cs675/f26/four-units.html}\par}
\end{frame}

% ---------------------------------------------------------------- 3
\begin{frame}[t]{Nearest neighbours, as attention}
  \vspace{1mm}
  \begin{columns}[c]
    \begin{column}{0.54\textwidth}
      \centering
      \figalt{A scatter of eight remembered Newark days, daily low against daily high, with a dashed vertical line at the query high of \AttQuery\ degrees and a red diamond at the predicted low of \AttYhat. Under the axis, a bar at each remembered day shows its attention weight; the bars near the query are tall and the far ones vanish.}
      \begin{tikzpicture}[x=1mm,y=1mm,line width=0.4pt]
        % highs 25..100 over 72 mm; lows 10..80 over 34 mm; a strip of weight bars underneath
        \pgfmathsetmacro{\xzero}{25}\pgfmathsetmacro{\xscale}{72/75}
        \pgfmathsetmacro{\yzero}{10}\pgfmathsetmacro{\yscale}{34/70}\pgfmathsetmacro{\ybase}{20}
        % the plot frame and its ticks
        \draw[faint] (0,\ybase) rectangle (72,{\ybase+34});
        \foreach \t in {30,50,70,90} { \draw[faint] ({(\t-\xzero)*\xscale},\ybase) -- ++(0,-1);
          \node[font=\tiny,color=muted,anchor=north] at ({(\t-\xzero)*\xscale},{\ybase-1.2}) {\t}; }
        \node[font=\tiny,color=muted,anchor=north] at (36,{\ybase-4.2}) {daily high ($^\circ$F)};
        \foreach \t in {20,40,60,80} { \node[font=\tiny,color=muted,anchor=east] at (-1,{\ybase+(\t-\yzero)*\yscale}) {\t}; }
        \node[font=\tiny,color=muted,rotate=90,anchor=south] at (-6,{\ybase+17}) {daily low ($^\circ$F)};
        % the weights, as bars under the plot
        \draw[faint] (0,8) -- (72,8);
        \node[font=\tiny,color=njred,anchor=east] at (-1,9.5) {weight $\alpha_i$};
        \foreach \x/\y/\a/\d in \AttDays {
          \fill[njred!75] ({(\x-\xzero)*\xscale-1.3},8) rectangle ({(\x-\xzero)*\xscale+1.3},{8+\a*8});
          \pgfmathparse{\a>0.05 ? 1 : 0}
          \ifnum\pgfmathresult=1 \node[font=\tiny,color=njred,anchor=north] at ({(\x-\xzero)*\xscale},7.4) {\a}; \fi }
        % the query
        \draw[dashed,njred,line width=0.5pt] ({(\AttQuery-\xzero)*\xscale},8) -- ({(\AttQuery-\xzero)*\xscale},{\ybase+34});
        \node[font=\tiny,color=njred,anchor=south] at ({(\AttQuery-\xzero)*\xscale},{\ybase+34.3}) {new day: high \AttQuery};
        % the remembered days, labelled by their lows
        \foreach \x/\y/\a/\d in \AttDays {
          \fill[ink] ({(\x-\xzero)*\xscale},{\ybase+(\y-\yzero)*\yscale}) circle (0.9);
          \node[font=\tiny,color=muted,anchor=south] at ({(\x-\xzero)*\xscale},{\ybase+(\y-\yzero)*\yscale+1}) {\y}; }
        % the answer
        \node[diamond,fill=njred,draw=white,line width=0.4pt,inner sep=0pt,minimum size=2.6mm] at ({(\AttQuery-\xzero)*\xscale},{\ybase+(\AttYhat-\yzero)*\yscale}) {};
        \node[font=\tiny,color=njred,anchor=west] at ({(\AttQuery-\xzero)*\xscale+2},{\ybase+(\AttYhat-\yzero)*\yscale}) {$\hat y=\AttYhat$};
      \end{tikzpicture}
    \end{column}
    \begin{column}{0.44\textwidth}
      {\small A regression by memory: eight remembered days $(x_i, y_i)$ and a new day with high $x$.}
      \par\vspace{2mm}
      {\small\begin{tabular}{@{}l@{\quad}l@{}}
        score, from the distance & $s_i = -\dfrac{d(x, x_i)^2}{2\tau^2}$ \\[3mm]
        \alert{softmax}: weights & $\alpha_i = \dfrac{e^{s_i}}{\sum_j e^{s_j}}$ \\[3mm]
        answer & $\hat y = \sum_i \alpha_i\, y_i$
      \end{tabular}}
      \par\vspace{2mm}
      {\footnotesize
      $k$-NN regression: $\alpha_i = 1/k$ for the $k$ nearest days and $0$ for the rest.
      Attention: every remembered day votes, with a weight that falls with its distance; $\tau$ plays the role of $k$.}
      \par\vspace{1.5mm}
      {\scriptsize\muteds{Here $\tau = \AttTau$. The three nearest days average \AttKnn; the least-squares line says \AttLine. Notes: 2.4, 2.12, 16.2.}}
    \end{column}
  \end{columns}
\end{frame}

\end{document}
